% Exact numerical abstract from the verified revision.
\begin{abstract}
Stopping quantum key generation need not immediately exhaust a trusted-relay network's stored keys. We study how finite reserves, interruption duration and fallback policy jointly determine timely key service and assurance loss. Generation, transport and online key-establishment failures are separated, with paired policy comparisons within each condition. A 360-run, five-topology grid includes 24-hour generation cuts with only 2\% of usable capacity initially filled. In the eight-node formation, fail-closed service is 6.70\%, compared with 97.53\% under fallback; compliant fallback service is only 6.61\%. The fraction of offered requests entitled to information-theoretically secure (ITS) key delivery but served below it is 90.08\% (95\% interval 89.20--90.97\%). With full starting reserves, the same cut produces no observed forced loss in the 24-hour cohort. An interpretable controller trained across several conditions raises mean count-based utility by 0.0011 on a held-out benchmark, yet increases associated below-ITS payload exposure by 1.71 percentage points while respecting every class exception cap. The results characterize modeled key service and contractual exposure, rather than payload throughput or field-validated security.
\end{abstract}
